\documentclass[11pt,a4paper]{article}
\usepackage[a4paper,margin=2.5cm]{geometry}
\setlength{\parindent}{0pt}
\setlength{\parskip}{0.6em}
\usepackage{helvet}
\renewcommand{\familydefault}{\sfdefault}
\usepackage[utf8]{inputenc}
\usepackage[T1]{fontenc}
\usepackage{setspace}
\onehalfspacing
\usepackage{titlesec}
\titleformat{\section}{\bfseries\Large\centering}{}{0pt}{}
\titlespacing*{\section}{0pt}{2em}{1em}

\begin{document}

\begin{center}
\textbf{\Large POLÍTICAS INSTITUCIONAIS}\\
\textbf{\large MOVIMENTO EM DEFESA DA EDUCAÇÃO -- MoveEduca}
\end{center}

\section*{1. FINALIDADE E APLICAÇÃO}

\textbf{1.1.} Estas Políticas Institucionais orientam a conduta, a governança, a integridade, a transparência e a atuação operacional do MoveEduca.

\textbf{1.2.} Aplicam-se a associados, dirigentes, conselheiros, empregados, voluntários, prestadores de serviços, parceiros e demais pessoas que atuem em nome ou em benefício da associação.

\textbf{1.3.} As políticas serão interpretadas de acordo com o Estatuto Social, o Regimento Interno, a legislação aplicável e os valores institucionais de bem do todo, união, neutralidade política, respeito, colaboração, integridade, sustentabilidade, inovação e transparência.

\section*{2. POLÍTICA DE INTEGRIDADE E CONDUTA}

\textbf{2.1.} A atuação institucional deverá observar legalidade, boa-fé, honestidade, respeito, responsabilidade, impessoalidade, transparência e compromisso com a educação.

\textbf{2.2.} É vedada qualquer prática de fraude, corrupção, favorecimento indevido, assédio, discriminação, retaliação, desvio de recursos, uso indevido da marca ou obtenção de vantagem pessoal incompatível com as finalidades do MoveEduca.

\textbf{2.3.} O MoveEduca manterá neutralidade político-partidária, sem apoio institucional a candidatos, partidos ou campanhas eleitorais, preservado o debate público qualificado sobre educação, cidadania e políticas públicas.

\textbf{2.4.} Violações a estas políticas deverão ser comunicadas à Diretoria Executiva, ao Conselho Fiscal ou ao Conselho de Administração, conforme a natureza do fato.

\section*{3. POLÍTICA DE CONFLITO DE INTERESSES}

\textbf{3.1.} Toda pessoa que participe de decisão institucional deverá declarar conflito de interesses real, potencial ou aparente.

\textbf{3.2.} A pessoa em conflito não poderá votar, aprovar, fiscalizar isoladamente ou influenciar a decisão correspondente.

\textbf{3.3.} Contratações com partes relacionadas somente serão admitidas com justificativa, transparência, compatibilidade de mercado, aprovação do órgão competente e registro documental.

\textbf{3.4.} A declaração de conflito de interesses deverá ser formalizada em termo próprio sempre que envolver contratação, parceria, remuneração, doação condicionada, seleção de fornecedor ou deliberação relevante.

\section*{4. POLÍTICA DE TRANSPARÊNCIA E PRESTAÇÃO DE CONTAS}

\textbf{4.1.} O MoveEduca prestará contas de suas atividades e recursos aos órgãos internos competentes, aos parceiros, aos financiadores e à Administração Pública quando aplicável.

\textbf{4.2.} A transparência deverá respeitar a proteção de dados pessoais, o sigilo legal, a confidencialidade estratégica e a segurança institucional.

\textbf{4.3.} Relatórios de atividades, demonstrações financeiras, pareceres e prestações de contas serão organizados conforme o Estatuto Social, as Normas Brasileiras de Contabilidade e as exigências dos instrumentos de parceria.

\section*{5. POLÍTICA DE PROTEÇÃO DE DADOS E CONFIDENCIALIDADE}

\textbf{5.1.} Dados pessoais serão coletados e utilizados apenas para finalidades legítimas relacionadas às atividades do MoveEduca.

\textbf{5.2.} O acesso a dados pessoais e informações internas será limitado às pessoas que necessitem dessas informações para o exercício de suas funções.

\textbf{5.3.} Informações de associados, beneficiários, crianças, adolescentes, parceiros, doadores e colaboradores deverão ser tratadas com cuidado, confidencialidade e observância da legislação aplicável.

\textbf{5.4.} Incidentes, perda de documentos ou exposição indevida de dados deverão ser comunicados imediatamente à Diretoria Executiva.

\section*{6. POLÍTICA DE VOLUNTARIADO}

\textbf{6.1.} A atuação voluntária deverá estar alinhada às finalidades do MoveEduca e poderá ser formalizada por termo próprio.

\textbf{6.2.} O voluntário deverá receber orientação mínima sobre a atividade, os responsáveis, a duração, os cuidados de conduta, a confidencialidade e a ausência de vínculo empregatício.

\textbf{6.3.} O ressarcimento de despesas de voluntários somente ocorrerá quando previamente autorizado e comprovado.

\section*{7. POLÍTICA DE COMPRAS E CONTRATAÇÕES}

\textbf{7.1.} Compras e contratações deverão observar necessidade real, finalidade institucional, economicidade, compatibilidade de mercado e documentação adequada.

\textbf{7.2.} Sempre que possível, a seleção de fornecedores deverá considerar mais de uma referência de preço, salvo urgência, exclusividade, baixo valor, continuidade técnica ou justificativa formal.

\textbf{7.3.} Contratações com recursos públicos seguirão o edital, o plano de trabalho, o instrumento firmado e a legislação aplicável.

\textbf{7.4.} É vedado fracionar despesas para evitar regras de aprovação, cotação, prestação de contas ou alçada.

\section*{8. POLÍTICA DE COMUNICAÇÃO E USO DE MARCA}

\textbf{8.1.} O uso do nome Movimento em Defesa da Educação, da denominação MoveEduca, do nome fantasia Move\&Educa e de marcas ou materiais institucionais dependerá de autorização da Diretoria Executiva ou de pessoa designada.

\textbf{8.2.} Comunicações públicas deverão ser verdadeiras, respeitosas, não discriminatórias e coerentes com a missão institucional.

\textbf{8.3.} É vedado manifestar opinião pessoal como se fosse posição oficial do MoveEduca sem autorização.

\section*{9. POLÍTICA DE GESTÃO DOCUMENTAL}

\textbf{9.1.} Documentos institucionais, financeiros, contábeis, fiscais, trabalhistas, societários, de projetos e de prestação de contas deverão ser arquivados de forma organizada.

\textbf{9.2.} A guarda documental observará os prazos legais, as exigências de parceiros e financiadores e a necessidade de comprovação de atos de gestão.

\textbf{9.3.} A eliminação de documentos relevantes dependerá de autorização da Diretoria Executiva, quando expirados os prazos legais e inexistente obrigação de guarda.

\section*{10. DISPOSIÇÕES FINAIS}

\textbf{10.1.} Estas Políticas Institucionais poderão ser detalhadas por manuais, formulários, modelos de termo, fluxos internos e deliberações dos órgãos competentes.

\textbf{10.2.} Os casos omissos serão resolvidos pela Diretoria Executiva, observados o Estatuto Social, o Regimento Interno e a legislação aplicável.

\vspace{2cm}
\begin{center}
\textbf{[CIDADE/UF]}, \textbf{[DATA]}.

\vspace{2cm}
\rule{8cm}{0.4pt}\\
\textbf{[NOME COMPLETO]}\\
Presidente do MoveEduca

\vspace{1.5cm}
\rule{8cm}{0.4pt}\\
\textbf{[NOME COMPLETO]}\\
Secretário(a)-Geral do MoveEduca

\vspace{1.5cm}
\rule{8cm}{0.4pt}\\
\textbf{[NOME DO ADVOGADO]}\\
\textbf{[OAB/UF Nº]}
\end{center}

\end{document}
